% Reconstructed proof module, 12 September 2026.
% This is newly authored mathematics, not a recovered src/energy.tex download.
% Required packages: amsmath, amssymb, amsthm; the standalone wrapper supplies them.
\section{Exact character-resolved collision minima and integral ES return}
\label{sec:reconstructed-character-energy}

\paragraph{Provenance and scope.}
This module reconstructs the character-resolved energy result stated in the
supplied joint ES discussion, \texttt{pasted-text.txt}, in its final
``Today 8:19 AM'' entry, subsection ``JT's contribution supplies an exact
refinement of the positivity test.'' The discussion credits the quadratic
character refinement to its integration of JT's bounded-divisor synthesis.
The original \texttt{src/energy.tex} was not available: the downloaded loose
\texttt{workbench.tex} is a wrapper referring to that missing file.
The shell and ordered channel conventions are those in the supplied earlier
\texttt{es\_dual\_frontier/workbench.tex} and the definition block in
\texttt{dual\_es.py}; their required identities are proved here in full.
The independently authored checker imports no downloaded code.
The result is a sufficient test on a specified full exponent box. It proves
no universal statement that every prime has a shell passing this test.

\subsection{The complete shell, exponent box, and ordered return}
Fix a prime $p=4m+1$, where $m\geq1$, and an integer
\[
 m+1\leq a\leq3m,\qquad R=4a-p.
\]
Then $3\leq R\leq8m-1=2p-3$, $R\equiv3\pmod4$, and $0<a<p$.
Since $p$ is prime, $\gcd(a,p)=1$; consequently
\[
 \gcd(a,R)=\gcd(a,p)=1,\qquad
 \gcd(p,R)=\gcd(p,4a)=1.
\]
Write the complete prime factorization as $a=\prod_{q\mid a}q^{e_q}$.
All these primes are units modulo $R$. Define
\[
 B_a=\prod_{q\mid a}\bigl([-e_q,e_q]\cap\mathbb Z\bigr),
 \quad v(\beta)=\prod_{q\mid a}q^{\beta_q}\in
 G=(\mathbb Z/R\mathbb Z)^\times,
 \quad u(\beta)=\prod_{q\mid a}q^{e_q+\beta_q}.
\]
Negative powers in $v$ mean multiplicative inverses in $G$.
Unique prime factorization proves that $\beta\mapsto u(\beta)$ is a bijection
from $B_a$ onto the positive divisors of $a^2$: its inverse is
$\beta_q=v_q(u)-e_q$. Multiplying the residue products gives
\[
 v(\beta)=u(\beta)a^{-1}\pmod R,
 \qquad u(-\beta)=a^2/u(\beta),
 \qquad v(-\beta)=v(\beta)^{-1}.
\]
In particular no unbounded subgroup or shorter exponent interval has
replaced $B_a$.

For any positive divisor $u\mid a^2$, put $d=a^2/u$ and retain the two
ordered rational triples
\begin{align*}
 Q_E(p,a,u)&=\left(a,\frac{pa+d}{R},\frac{pa+p^2u}{R}\right),\\
 Q_M(p,a,u)&=\left(a,\frac{p(a+d)}{R},\frac{p(a+u)}{R}\right).
\end{align*}
Their first denominator remains distinguished; the last two entries are
not sorted. In either case, if these entries are denoted $Y,Z$, then
\[
 (RY-pa)(RZ-pa)=p^2a^2.
\]
Indeed the two factors are $(d,p^2u)$ in channel $E$ and $(pd,pu)$ in
channel $M$. Expanding and cancelling the identical constant term gives
$R^2YZ-paR(Y+Z)=0$. Since $R,p,a,Y,Z$ are positive, division gives
\[
 \frac1Y+\frac1Z=\frac{R}{pa},\qquad
 \frac1a+\frac1Y+\frac1Z=\frac{p+R}{pa}=\frac4p.
\]
Both triples therefore always give positive rational ES decompositions.

The exact integral gates follow without discarding the modulus $R$.
For $E$, multiplication of the first numerator by the unit $u$ yields
\[
 u(pa+d)\equiv a^2(4u+1)\pmod R,
 \qquad pa+p^2u\equiv4a^2(4u+1)\pmod R.
\]
For $M$ the exact numerator identities are
\[
 u\,p(a+d)=pa(u+a),\qquad p(a+u)=p(u+a).
\]
Since every multiplier used is a unit modulo $R$, the two reduced
denominators in each channel coincide and are respectively
\[
 D_E=\frac{R}{\gcd(R,4u+1)},\qquad
 D_M=\frac{R}{\gcd(R,u+a)}.
\]
Thus $Q_E$ is integral exactly when $R\mid4u+1$, and $Q_M$ exactly when
$R\mid u+a$. This proves all assertions about integrality, not merely
one implication.

Since $u=av(\beta)$ modulo $R$ and $4a=p$ modulo $R$, these tests become
\[
 M:\ v(\beta)=-1,
 \qquad E:\ v(\beta)=-p^{-1}.
\]
The complete box is stable under $\beta\mapsto-\beta$, and
$(-p)^{-1}=-p^{-1}$, so a vector with residue $-p$ also supplies an
exterior witness by applying this involution. Consequently simultaneous
failure of both channels in this shell is equivalent to absence of all
the residues in the deduplicated, inversion-stable set
\[
 T=\{-1,-p,-p^{-1}\}\subseteq G.
\]
The set may have fewer than three elements, and may contain the identity;
both possibilities are retained below.

For the additional integral marking put
\[
 g=\gcd(a,u),\quad h=g^2/u,\quad r=u/g,\quad s=a/g.
\]
If $e=v_q(a)$ and $f=v_q(u)$, then $0\leq f\leq2e$ and the exponents
of $h,r,s$ are $2\min(e,f)-f$, $f-\min(e,f)$ and
$e-\min(e,f)$. The first is $f$ when $f\leq e$ and $2e-f$ otherwise,
so all are nonnegative. At least one of the last two is zero, proving
$\gcd(r,s)=1$. Direct substitution now proves
\[
 a=hrs,\quad u=hr^2,\quad d=hs^2,
\]
and the ordered triples are
\[
 Q_E=(hrs,hs\kappa,phr\kappa),\quad\kappa=(pr+s)/R,
 \qquad
 Q_M=(hrs,phs\lambda,phr\lambda),\quad\lambda=(r+s)/R.
\]
Since $g$ is a unit modulo $R$, the middle gate is equivalent to
$R\mid r+s$. The exterior numerator satisfies
$pa+d=hs(pr+s)$, where $hs$ is a unit modulo $R$, so the exterior
gate is equivalent to $R\mid pr+s$. Hence these quotients are positive
integers whenever the corresponding target occurs. Every marking returns
to the original objects via $a=hrs$, $u=hr^2$ and
$\beta_q=v_q(u)-e_q$; moreover
\[
 u=\frac{a^2}{RY-pa}\quad(E),\qquad
 u=\frac{pa^2}{RY-pa}\quad(M).
\]
Thus the channel and ordered denominators retain an explicit inverse.

For completeness, the full interval of shells suffices for the ES
existence question at this $p$. Given any positive integral solution,
move a smallest denominator to the distinguished position $a$, retaining
that permutation. The inequalities
$4/p>1/a$ and $4/p\leq3/a$ give $p/4<a\leq3p/4$, hence
$m+1\leq a\leq3m$. Its other two denominators obey the residual product
identity above, and
$RY-pa=paY/Z>0$, $RZ-pa=paZ/Y>0$. Since $p\nmid a$, the first residual
is uniquely $p^i v$ with $i\in\{0,1,2\}$ and $v\mid a^2$.
For $i=0$, take $E,u=a^2/v$ in the existing order. For $i=1$, take
$M,u=a^2/v$ in the existing order. For $i=2$, take $E,u=v$ and retain
the swap of the last two entries. The displayed residuals verify each
case. Conversely, an integral candidate already solves ES. This proves
the exact equivalence with a hit in the full atlas; it does not assert
that such a hit always exists.

\subsection{Multiplicities, parity, and character cells}
Define
\[
 \mu(g)=\#\{\beta\in B_a:v(\beta)=g\},\quad
 N=\sum_{g\in G}\mu(g)=\prod_{q\mid a}(2e_q+1),\quad
 E=\sum_{g\in G}\mu(g)^2.
\]
The involution $\beta\mapsto-\beta$ bijects the fibres of $g$ and
$g^{-1}$, proving $\mu(g)=\mu(g^{-1})$. Its only fixed vector in the
integer box is zero. The fibre over $1$ contains that vector and otherwise
splits into two-element orbits; thus $\mu(1)$ is odd and at least one.
If $g^2=1$ but $g\ne1$, its fibre is stable under this involution and
contains no fixed vector, so $\mu(g)$ is even.

The collision count also has the independent exact expression
\[
 E=\sum_{\substack{-2e_q\leq\delta_q\leq2e_q\ (q\mid a)\\
                  \prod q^{\delta_q}=1\ (\bmod R)}}
       \prod_{q\mid a}(2e_q+1-|\delta_q|).
\]
Indeed $E$ counts ordered pairs $(\beta,\gamma)$ with identical residue.
Their coordinate difference $\delta=\beta-\gamma$ has residue one.
For a single interval $[-e,e]$, the number of ordered pairs of difference
$\delta$ is $2e+1-|\delta|$ for $|\delta|\leq2e$, and zero otherwise:
the first entry lies in $[-e,e]\cap[-e+\delta,e+\delta]$.
Taking the product of these counts proves the formula.

Let $V$ be a finite elementary abelian $2$-group, written
multiplicatively, and let $\pi:G\to V$ be a homomorphism; surjectivity
is not required. For $v\in V$ put
\[
 C_v=\pi^{-1}(v),\qquad
 N_v=\sum_{g\in C_v}\mu(g),\qquad
 E_v=\sum_{g\in C_v}\mu(g)^2.
\]
As $v^{-1}=v$, each cell is stable under inversion. Its mass is even
except in the cell containing $1$, where its mass is odd: inverse pairs
contribute twice their common multiplicity, nonidentity involutions have
even multiplicity, and the identity contributes odd multiplicity.
Empty cells have mass and energy zero.

For each odd prime $\ell\mid R$, the prime-local quadratic character
$\chi_\ell(g)$ is $1$ when $g$ is a nonzero square modulo $\ell$ and
$-1$ otherwise. This is a homomorphism: in the multiplicative group of
the field $\mathbb Z/\ell\mathbb Z$, the square map is a homomorphism
whose kernel is exactly $\{1,-1\}$, since
$(x-1)(x+1)=0$ implies $x=1$ or $x=-1$ in a field. Every fibre has
two elements, so the square subgroup has index two. The quotient by
that subgroup has two elements and therefore supplies precisely this
sign character. Reduction modulo $\ell$ is a group homomorphism on $G$,
which proves the assertion for $\chi_\ell$ on $G$ as well.

Writing $R=\prod_{\ell\mid R}\ell^{f_\ell}$, define
\[
 \pi_{\rm quad}(g)=(\chi_\ell(g))_{\ell\mid R},\qquad
 \pi_{\rm Jac}(g)=\prod_{\ell\mid R}\chi_\ell(g)^{f_\ell}.
\]
Both are homomorphisms. The second is obtained from the first by the
explicit homomorphism $(\epsilon_\ell)\mapsto
\prod\epsilon_\ell^{f_\ell}$; factors with even $f_\ell$ disappear only
from this product. For example, modulo $9$ the residue $2$ is a
non-square modulo $3$, so its prime-local character is $-1$, whereas
its Jacobi character is $(-1)^2=1$. Thus these two partitions need
not coincide. The trivial homomorphism is the unpartitioned case.

\subsection{The exact minimum, including all impossible cases}
For a fixed cell let $A_v=C_v\setminus T$. Since both $C_v$ and $T$
are inversion-stable, so is $A_v$. Its slots are the unordered inverse
pairs $\{g,g^{-1}\}$ with $g\ne g^{-1}$, the individual nonidentity
involutions $g^2=1$, $g\ne1$, and the identity if it is present.
Write $b_v=1$ when $1\in C_v$ and $b_v=0$ otherwise.

For a nonnegative integer mass $n$, define $L_v(n)$ as the minimum of
$\sum_{g\in C_v}\nu(g)^2$ over all maps
$\nu:C_v\to\mathbb Z_{\geq0}$ satisfying
\[
 \sum_{g\in C_v}\nu(g)=n,\quad
 \nu(g)=\nu(g^{-1}),\quad
 \nu|_{T\cap C_v}=0,
\]
and additionally $\nu(1)$ odd when $1\in C_v$, and $\nu(g)$ even
for every nonidentity involution in the cell. If there is no such map,
set $L_v(n)=+\infty$.

Here is a complete evaluation. If $1\in T\cap C_v$, the minimum is
$+\infty$ for every $n$, since an odd nonnegative identity value cannot
vanish. Otherwise the identity slot exists precisely when $b_v=1$.
If $n<b_v$ or $n-b_v$ is odd, the minimum is $+\infty$.
In all remaining cases put $K=(n-b_v)/2$. Write $P_v$ for the set
of inverse-pair slots and $J_v$ for the nonidentity involution slots.
Use nonnegative integer variables $k_P$ for $P\in P_v$, $z_j$ for
$j\in J_v$, and $k_0$ only if $b_v=1$. Then
\[
 L_v(n)=\min_{\sum_P k_P+\sum_j z_j+b_vk_0=K}
 \left(2\sum_P k_P^2+4\sum_j z_j^2+
 \begin{cases}(1+2k_0)^2,&b_v=1,\\0,&b_v=0.\end{cases}\right),
\]
where the absent variable is omitted from both sums. If $K>0$ and
there are no slots, the minimum is $+\infty$; if $K=0$ it is $b_v$.
These are all the impossible cases: when at least one slot exists, it
can hold every remaining pair of mass units. For a pair slot the two
values are $(k_P,k_P)$, for a nonidentity involution its value is
$2z_j$, and for the identity its value is $1+2k_0$. These assignments
give a bijection between feasible slot allocations and feasible maps
$\nu$, proving the optimization formula. There are finitely many
allocations for fixed $K$, so every finite minimum is attained.

Start the identity at mass one and energy one when $b_v=1$ and start
every other slot at zero. Incrementing a slot with current index $k$
by one adds two units of mass. Its exact energy increment is respectively
\[
 \begin{array}{c|c|c}
 \text{slot}&\text{current values}&\text{incremental energy}\\\hline
 \{g,g^{-1}\},\ g\ne g^{-1}&(k,k)&2(k+1)^2-2k^2=4k+2\\
 g^2=1,\ g\ne1&2k&(2k+2)^2-(2k)^2=8k+4\\
 g=1&1+2k&(3+2k)^2-(1+2k)^2=8k+8
 \end{array}
\]
In each row $k$ starts at zero. Thus the marginal sequences are
$2,6,10,\ldots$, $4,12,20,\ldots$, and $8,16,24,\ldots$.
Selecting the $K$ least costs from the union of these increasing
sequences proves the minimum exactly. To justify feasibility, every
predecessor in a sequence is strictly smaller than its successor,
so selecting a successor among the $K$ least costs selects its whole
prefix. Ties between different sequences can be resolved in any fixed
order. To justify optimality, any feasible allocation selects exactly
$K$ costs from this union, hence its sum is at least the sum of the
$K$ least ones. Equivalently, replacing a selected unnecessarily large
cost by the least unselected available cost improves the allocation
until the same minimum is reached. A heap holding the first unused
cost of each sequence implements this finite procedure.

\paragraph{Forcing theorem.}
For the actual full-box multiplicities and every cell $v$,
\[
 \mu|_{T\cap C_v}=0\quad\Longrightarrow\quad E_v\geq L_v(N_v).
\]
Indeed all other constraints defining $L_v$ were proved for $\mu$,
so this restriction is a feasible map if its target fibres vanish.
Taking the contrapositive proves
\[
 E_v<L_v(N_v)\quad\Longrightarrow\quad
 \text{there is }\beta\in B_a\text{ with }v(\beta)\in T\cap C_v.
\]
The preceding explicit map then produces a positive integral ES
witness. In particular, if $T\cap C_v=\{-1\}$, the witness is in
the middle channel. A finite energy is smaller than $+\infty$,
so incompatible cells are included without an extra convention.
If $1\in T$, the zero vector itself is already a target and is
detected by this rule. If $T\cap C_v$ is empty, the actual restriction
is feasible and this cell can never produce a false certificate.

Every finite $L_v(n)$ is attained by a distribution satisfying exactly
the specified mass, symmetry, parity, and exclusions. Hence it is the
largest universal lower bound on energy from those constraints alone.
This is the precise optimality assertion. It does not state that every
energy above the minimum is attainable, or that every attaining abstract
distribution arises from the arithmetic exponent box.

\subsection{Exact refinement and preservation of certificates}
Let $\pi_2:G\to V_2$ and $\pi_1:G\to V_1$ be such homomorphisms,
and suppose $\pi_1=\rho\circ\pi_2$ for a homomorphism
$\rho:V_2\to V_1$. A coarse cell $C_v$ is the disjoint union of
the fine cells $C_w$ with $\rho(w)=v$. For any prescribed fine
masses $n_w$, with $n=\sum_{\rho(w)=v}n_w$, the exact inequality is
\[
 L^{(1)}_v(n)\leq\sum_{\rho(w)=v} L^{(2)}_w(n_w).
\]
All terms take values in $\mathbb Z_{\geq0}\cup\{+\infty\}$.
If some right-hand term is infinite the inequality is automatic.
Otherwise choose a minimizer on each fine cell and join these maps.
Each inversion orbit lies entirely in a fine cell, so the joined map
obeys every coarse constraint, with mass $n$ and energy equal to the
sum on the right. This proves the inequality. Restricting any coarse
feasible map to its fine cells proves the stronger exact identity
\[
 L^{(1)}_v(n)=
 \min_{\substack{n_w\in\mathbb Z_{\geq0}\\
                  \sum_{\rho(w)=v}n_w=n}}
       \sum_{\rho(w)=v}L^{(2)}_w(n_w).
\]
If the coarse minimum is finite, its minimizer supplies the masses
and gives the reverse inequality. If it is infinite, no combination
of finite fine minimizers can exist, since their join would be a
coarse feasible map. This also proves equality in the infinite case.

For actual masses and energies,
$N^{(1)}_v=\sum N^{(2)}_w$ and $E^{(1)}_v=\sum E^{(2)}_w$ by disjointness.
If every fine cell satisfies $E^{(2)}_w\geq L^{(2)}_w(N^{(2)}_w)$,
their sum and the inequality above give
$E^{(1)}_v\geq L^{(1)}_v(N^{(1)}_v)$. The contrapositive proves that
every coarse forcing certificate yields a forcing certificate in at
least one of its fine cells. This applies successively to the trivial,
Jacobi, and all-prime-local quadratic partitions, with the explicit
product map already proved above.

\subsection{The complete numerical witness at p = 1201}
Take $p=1201$, $a=310=2\cdot5\cdot31$, $R=39=3\cdot13$.
To certify primality explicitly, $34^2<1201<35^2$ and the residues
of $1201$ modulo the primes $2,3,5,7,11,13,17,19,23,29,31$ are
respectively $1,1,1,4,2,5,11,4,5,12,23$. None is zero, and every
composite integer has a prime divisor at most its square root, as
follows by writing it as the product of two positive factors larger
than one and taking a prime factor of the smaller one. These are
all primes at most 34, so $1201$ is prime.
There are exactly $3^3=27$ exponent vectors with each coordinate
in $\{-1,0,1\}$. In the residue ring, $p=31$ and $31^{-1}=34$,
since $31\cdot34=27\cdot39+1$. Hence $T=\{5,8,38\}$.
The complete histogram on the 24 units is
\[
\begin{array}{c|rrrrrrrrrrrr}
g&1&2&4&5&7&8&10&11&14&16&17&19\\\hline
\mu(g)&1&1&1&1&2&1&1&0&2&1&1&2\\[2pt]
g&20&22&23&25&28&29&31&32&34&35&37&38\\\hline
\mu(g)&1&1&1&0&2&1&1&0&1&1&2&2
\end{array}
\]
This table can be checked by multiplying all 27 combinations of
$(2^{-1},1,2)=(20,1,2)$, $(5^{-1},1,5)=(8,1,5)$ and
$(31^{-1},1,31)=(34,1,31)$ modulo $39$. Its six entries of value
two and fifteen entries of value one give $N=27$ and $E=6\cdot4+15=39$.

Deleting $T$ leaves nine inverse pairs,
\[
 \{2,20\},\{4,10\},\{7,28\},\{11,32\},\{16,22\},
 \{17,23\},\{19,37\},\{29,35\},\{31,34\},
\]
the two nonidentity involutions $14,25$, and the identity. The
unpartitioned mass budget is $(27-1)/2=13$. The 13 least marginal
costs are nine copies of $2$, two copies of $4$, and two copies of
$6$. Thus
\[
 L_{\mathrm{unpartitioned}}(27)=1+9\cdot2+2\cdot4+2\cdot6=39.
\]
Equality $E=39$ gives no strict forcing certificate at this level.

The negative Jacobi cell is exactly
\[
 C_- =\{7,14,17,19,23,28,29,31,34,35,37,38\}.
\]
For a direct sign check the nonzero squares modulo $3$ are $\{1\}$,
and those modulo $13$ are $\{1,3,4,9,10,12\}$; taking the product
of these two signs gives exactly the displayed cell. Its table entries
give $N_-=18$, $E_-=30$, and $T\cap C_-=\{38\}$.
Removing this target leaves five pairs
\[
 \{7,28\},\{17,23\},\{19,37\},\{29,35\},\{31,34\}
\]
and the nonidentity involution $14$. There is no identity slot.
The budget is nine, whose nine least costs are five copies of $2$,
one copy of $4$, and three copies of $6$. Therefore
\[
 L_-(18)=5\cdot2+4+3\cdot6=32
 =\min_{n_1+\cdots+n_5+z=9}
       \left(2\sum_{i=1}^5 n_i^2+4z^2\right).
\]
All variables in the minimum are nonnegative integers. The allocation
$(n_1,\ldots,n_5,z)=(2,2,2,1,1,1)$ attains 32, and the marginal
argument proves that no smaller value occurs. Hence $30<32$ forces
the middle target in this specific cell.

The finer labels are ordered by the primes $(3,13)$; their complete
mass, energy, minimum and target data are
\[
\begin{array}{c|r|r|r|l}
v&N_v&E_v&L_v(N_v)&T\cap C_v\\\hline
(-1,-1)&4&4&4&\{5,8\}\\
(-1,+1)&8&12&14&\{38\}\\
(+1,-1)&10&18&18&\varnothing\\
(+1,+1)&5&5&5&\varnothing
\end{array}
\]
These values follow by grouping the full histogram using the explicit
square lists. The allowed slot counts are respectively two pairs;
two pairs and the involution $14$; three pairs; and two pairs, the
identity, and the involution $25$. The marginal formula gives the
four displayed minima, including $2+2+4+6=14$ for $(-1,+1)$.

An actual middle fibre vector is $\beta=(0,-1,-1)$ in the prime
order $(2,5,31)$. Its residue is $8\cdot34=272=6\cdot39+38$, and
its original divisor is $u=2^{1+0}5^{1-1}31^{1-1}=2$.
The complete marking is therefore
\[
 (p,R,u,h,r,s,\lambda)=(1201,39,2,2,1,155,4),
\]
because $\gcd(310,2)=2$ and $(1+155)/39=4$. Substitution in the
ordered middle return gives
\[
 Q_M=(310,\ 1201\cdot2\cdot155\cdot4,\ 1201\cdot2\cdot1\cdot4)
     =(310,1489240,9608).
\]
Every number is a positive integer. The earlier proved reciprocal
identity now verifies the ES equation exactly, and the retained
slope satisfies $r/s=1/155<1/8$. This is a stronger forcing
certificate for this witness, with its entire divisor vector retained.

\subsection{Finite verification and the boundary of the result}
The independent standard-library checker compares the marginal algorithm
with an independent dynamic program
that uses the occupancy squares directly. It includes all masses
$0\leq n\leq19$, zero through three pair slots, zero through three
nonidentity involution slots, presence or absence of the identity,
and the forbidden-identity case. It also compares the actual cells
of the $1201,310$ example with that dynamic program. The checker
uses explicit exceptions, so optimized Python retains every check.
For its arithmetic cases it enumerates the full box and independently
enumerates divisors of $a^2$ by trial divisibility and complementary
divisors. The signed-box symmetries, parity, target exclusions,
refinement inequalities, and extracted rational identities are checked
exactly. The weighted autocorrelation is independently checked on
the listed examples, including shells with $R\geq p$.

The scan uses exactly the primes $p\equiv1\pmod{12}$ with $p\leq1500$
and the first-half shells $m+1\leq a\leq2m$, equivalently $R<p$.
It contains 54 primes and 9531 shells and gives respectively
214 unpartitioned, 231 Jacobi, and 240 fully quadratic certifications.
These are shell counts. The general theorem retains the full shell
range through $3m$; this finite scan does not replace that range.
The empty shell $p=37,a=18,R=35$ is retained and is not certified by
any partition. Ordinary and optimized execution reproduce the same
mathematical JSON outputs, with execution details recorded separately.
These finite results and the proved sufficient criterion do not
establish that every prime has a certified shell, and a shell that
does not pass this strict inequality may still have a genuine target.
